\section{Conclusion and Future Work}\label{sec:conclusion}
We compared four personalisation mechanisms, local-only training and DP-FedAvg under one record-level privacy
budget. Differential privacy hits the weakest clients hardest: at $\eps=2$ DP-FedAvg loses 13.1 pts of mean accuracy
but 18.3 pts on its worst-10\% clients. Personalisation more than reverses this. At $\eps=2$, FedPer raises mean
accuracy from 27.2\% to 54.1\% and worst-10\% accuracy from 12.1\% to 38.1\%, improves every client, and costs no more
communication than DP-FedAvg (22.4 MB per round). Local-only training reaches a slightly higher mean (56.2\%) with
no communication at all, which shows that under strong label skew the gain comes from fitting each client's own
label mix rather than from sharing. Below $\eps=1$ the ranking changes: fine-tuning and Per-FedAvg lose the least
accuracy per unit of privacy and keep the worst clients highest. The q-FFL fairness objective did not help; under DP
it stalled the global model. Clustering (IFCA) gave the smallest gain at twice the communication. In short, with
strong label skew, the most personalisation per unit of privacy comes from keeping a client-specific part of the
model---FedPer---at moderate budgets, and from light local adaptation of a shared model at small budgets.

Future work follows from the limitations above: client-level DP with thousands of clients per round on naturally
federated data such as FEMNIST; minimax fairness through agnostic FL~\cite{mohri2019agnostic} under the same
equal-budget protocol; q-FFL with a server step size retuned for the DP noise level; stronger likelihood-ratio membership
attacks~\cite{carlini2022membership}; modern census
data such as ACS Income~\cite{ding2021retiring} for group fairness; per-client privacy budgets for clients with very
little data; and the effect of wireless channel constraints---partial participation and limited uplink---on which
mechanism is most efficient.
